% \documentclass[../main.tex]{subfiles}
% \begin{document}

\section{Introduction}
% Kind of vacuous but okay.
\IEEEPARstart{C}{omputer} science has become increasingly integrated into modern society.
% Maybe a citation? Maybe vacuous? Introducing the concept at hand is _hard_.
Consequentially, educational systems have placed greater emphasis on preparing students to understand and engage with computational technologies.
% Citation?
Secondary computer science education has expanded substantially over the past two decades, driven by workforce demands, technological innovation, and
  growing recognition that computational thinking represents an essential component of contemporary literacy.
% This one might need a citation, too.
Within this broader movement, algorithms occupy a central role because they provide the procedural foundations through which computational problems
  are analyzed, represented, and solved.

% Citation?
Despite their importance, algorithms remain among the more challenging concepts for secondary students to learn and for educators to teach.
% Citation?
Algorithmic thinking requires learners to reason abstractly, decompose complex problems into manageable components, recognize patterns, and develop
  step-by-step procedures capable of producing desired outcomes.
% Citation?
These skills often differ from those emphasized in traditional academic subjects and may be complicated by student misconceptions, difficulties with
  abstraction, varying levels of prior programming experience, and challenges transferring knowledge between mathematics and computing contexts.
% Citation (this one should be easy)?
As computer science programs continue to expand, educators and researchers have increasingly sought to identify effective methods for teaching
  algorithms, assessing student understanding, and preparing teachers to support algorithmic learning.

% Good outline.
This literature review examines current research related to the teaching of algorithms in secondary education.
The review first establishes the educational context surrounding algorithms and the growth of computer science education.
It then explores instructional approaches used to teach algorithmic concepts, challenges students encounter when learning algorithms, methods for
  assessing algorithmic knowledge, and the role of teacher preparation and professional development in supporting effective instruction.
By synthesizing findings across these areas, this review aims to provide a comprehensive overview of the current state of research on algorithm
  education and identify considerations relevant to future instructional practice and educational research.

% \end{document}
